\section{Referencial teórico}

\subsection{Crescer testando hipóteses}

Em um experimento controlado com 116 startups italianas, Camuffo e coautores verificaram que empreendedores treinados para formular e testar hipóteses de forma rigorosa obtiveram melhor desempenho e pivotaram mais vezes, sem maior taxa de abandono [4]. O resultado sustenta o método da Lean Startup e do desenvolvimento de clientes, em que o negócio avança por ciclos de construir, medir e aprender [5, 6, 7]. Para o PetCuida, cada etapa do plano tem uma hipótese, um teste e um critério de avanço.

\subsection{Risco da escala prematura}

O relatório Startup Genome, baseado em dados de mais de 3.200 startups, apontou a escala prematura (contratar, gastar ou expandir antes de validar o produto) como a causa mais comum de fracasso entre as empresas analisadas [8]. Por isso, o plano só amplia território depois de atingir critérios mínimos de ativação, retenção e oferta local.

\subsection{Plataforma de dois lados}

O PetCuida conecta tutores e prestadores de serviço, o que o caracteriza como plataforma multilateral. Rochet e Tirole mostram que, nesses mercados, a estrutura de preços entre os lados é decisiva para atrair participantes [9]; Parker, Van Alstyne e Choudary destacam o problema do ``ovo e da galinha'': sem oferta, a demanda não fica; sem demanda, a oferta não entra [10]. A consequência prática é começar subsidiando o lado mais sensível a preço (o tutor) e garantir oferta mínima antes de divulgar.

\subsection{Matriz de crescimento e métricas}

A matriz de Ansoff classifica o crescimento em penetração, desenvolvimento de mercado, desenvolvimento de produto e diversificação [11]. As métricas do plano se apoiam em valor do cliente ao longo do tempo [12], na recomendação líquida (NPS) [13] e nos funis de ativação e retenção descritos em obras de crescimento e analítica de startups [14, 15, 16].
